\documentclass[11pt,a4paper]{article}
\usepackage[margin=25mm]{geometry}
\usepackage[T1]{fontenc}
\usepackage{lmodern,amsmath,amssymb,amsthm,microtype}
\usepackage[hidelinks]{hyperref}
\hypersetup{pdftitle={Reducing the map families in pruning}}
\newtheorem{proposition}{Proposition}
\newtheorem{lemma}[proposition]{Lemma}
\newcommand{\Hom}{\operatorname{Hom}_0}
\newcommand{\End}{\operatorname{End}_0}
\newcommand{\im}{\operatorname{im}}
\newcommand{\Span}{\operatorname{span}_{\mathbb F_2}}
\newcommand{\F}{\mathbb F_2}
\begin{document}

\section{Reducing the map families in pruning}
Let $M$ be a finitely presented bigraded persistence module over $\F$,
$N=M(\delta)$ for $\delta=(2\epsilon,2\epsilon)$ with $\epsilon\ge0$, and $\eta:M\to N$
the canonical structure map. Write $E=\End(M)$ and $H=\Hom(M,N)$,
where all maps have degree zero. The pruning passes are
\begin{align}
 \mathcal I(U)&=U\cap\bigcap_{f\in H}f^{-1}(\eta U),\label{eq:i}\\
 \mathcal K(V)&=\eta^{-1}\!\left(\eta V+\sum_{f\in H}f(V)\right),
 \label{eq:k}
\end{align}
iterated from $I_0=M$ and $K_0=0$. A vector-space basis of $H$ suffices
in both formulas. We reduce the number of image and preimage computations
further by removing redundant lifts, choosing generators under composition,
and discarding dependencies that arise on the current submodule.

\subsection{Actual map classes and inherited endomorphisms}
Fix a graded presentation $F_1\xrightarrow{R_M}F_0\twoheadrightarrow M$.
Let $L_\delta$ be the vector space of lifts $F_0\to F_0(\delta)$ that
send source relations into target relations, and $Z_\delta$ its subspace of lifts inducing zero
module maps. Thus $H=L_\delta/Z_\delta$ and $E=L_0/Z_0$.
Rather than computing redundant lifts first, actual Hom classes can be
obtained directly from the source generator grades $g_i$ and relation
grades $r_j$:
\begin{equation}\label{eq:hom}
 H\cong\ker\!\left(\bigoplus_i N_{g_i}
 \xrightarrow{\ D\ }\bigoplus_j N_{r_j}\right),\qquad
 D(v)_j=\sum_{i:(R_M)_{ij}\ne0}(R_M)_{ij}
 N(g_i\le r_j)(v_i).
\end{equation}
Indeed, a tuple $(v_i)$ uniquely specifies a map $F_0\to N$, and this
map factors through $M$ precisely when it annihilates $R_M$. Taking
each $N_g$ as a cokernel imposes the target relations before the kernel
computation. The same construction with target $M$ gives $E$.

The inherited maps $\eta e$, $e\in E$, can also be omitted whenever the
current submodule is stable under all endomorphisms. We therefore work in
\begin{equation}\label{eq:q}
 Q=H/(\eta E)\cong L_\delta/(Z_\delta+\tau L_0),
\end{equation}
where $\tau$ transports an endomorphism lift to the shifted target.
Naturality gives $e^\delta\eta=\eta e$, with $e^\delta=e(\delta)$.
Consequently, $Q$ has well-defined left and right $E$-actions
\[
 e\cdot[f]=[e^\delta f],\qquad [f]\cdot e=[fe],
\]
since $e^\delta(\eta a)=\eta(ea)$ and $(\eta a)e=\eta(ae)$.

To compute $Q$ using representative matrices, first span $Z_\delta$:
for each source generator, insert the target relation columns available
at its grade into that source column, with all other columns zero.
These span exactly the lifts whose generator values are zero in $N$.
Reduce the Hom representatives modulo this space and the transported
End representatives, retaining independent residues. This computes
\eqref{eq:q} independently of the chosen lifts, since
$\tau Z_0\subseteq Z_\delta$. Thus
$\dim Q=\dim H-\operatorname{rank}(e\mapsto\eta e)$; the rank need
not equal $\dim E$.

\subsection{Generators under composition}
Choose representatives $s_1,\ldots,s_\ell$ and $t_1,\ldots,t_r$ such that
\begin{equation}\label{eq:generators}
 Q=\sum_{i=1}^{\ell}E[s_i],\qquad
 Q=\sum_{j=1}^{r}[t_j]E.
\end{equation}
The I pass uses the postcomposition generators $s_i$; the K pass uses
the precomposition generators $t_j$. One selected map can generate
several independent classes under the action, so these families can be
smaller than a vector-space basis of $Q$.

\begin{lemma}\label{lem:invariant}
Every complete pass of~\eqref{eq:i}--\eqref{eq:k}, starting at $M$ and
$0$, yields submodules stable under $E$.
\end{lemma}
\begin{proof}
Suppose $U$ is stable under $E$ and $x\in\mathcal I(U)$. For $e\in E$,
we have $ex\in U$ and $f(ex)=(fe)(x)\in\eta U$ for every $f\in H$,
since $fe\in H$. Hence $ex\in\mathcal I(U)$.
For K, the sum $L=\eta V+\sum_{f\in H}f(V)=\sum_{f\in H}f(V)$
is stable under every $e^\delta$, since $e^\delta f\in H$.
If $\eta x\in L$, naturality implies $\eta(ex)=e^\delta\eta x\in L$.
Thus $\mathcal K(V)$ is stable under $E$. The initial submodules $M$
and $0$ are stable, completing the induction.
\end{proof}

\begin{proposition}\label{prop:families}
If $U,V$ are stable under $E$ and~\eqref{eq:generators} holds, then
\[
 \mathcal I(U)=U\cap\bigcap_i s_i^{-1}(\eta U),\qquad
 \mathcal K(V)=\eta^{-1}\!\left(\eta V+\sum_j t_j(V)\right).
\]
\end{proposition}
\begin{proof}
Every $f\in H$ can be written in both forms
\[
 f=\eta a+\sum_i b_i^\delta s_i,
 \qquad f=\eta a'+\sum_j t_jc_j,
 \qquad a,a',b_i,c_j\in E.
\]
For I, if $x\in U$ satisfies $s_i(x)\in\eta U$ for all $i$, then
$a(x)\in U$ and $b_i^\delta(\eta U)=\eta b_i(U)\subseteq\eta U$.
The first expression therefore gives $f(x)\in\eta U$ for every $f$.
The converse follows because $s_i\in H$.
For K, stability of $V$ and the second expression give
$f(V)\subseteq\eta V+\sum_j t_j(V)$. The reverse inclusion in the
full image sum follows from $\eta,t_j\in H$. Taking the preimage under
the same $\eta$ proves the second equality.
\end{proof}

The lemma and proposition show inductively that the smaller families
preserve every complete pass and hence the stabilized $I$, $K$, and
$I/(I\cap K)$. A vector-space basis of $Q$ also suffices, by taking
scalar multiples of the identity as coefficients. Stability is required
at the start of a pass; intermediate I intersections need not be stable.

\paragraph{Computing the generators.}
Let $q_1,\ldots,q_d$ represent a basis of $Q$, and let $e_1,\ldots,e_t$
be a basis of $E$. Tabulate the coordinates of $[e_a^\delta q_j]$ and
$[q_je_a]$ in $Q$. For either action, maintain the generated subspace
$W$, initially zero, and choose a candidate $s$ maximizing
$\dim(W+E[s])$ on the left or $\dim(W+[s]E)$ on the right.
Candidates include all $q_j$, their total sum,
and a fixed number of random linear combinations. Extend $W$ by the
chosen orbit and repeat until $\dim W=d$.

This procedure has an exact termination certificate. If $W\ne Q$, some
$[q_j]\notin W$; because the identity belongs to $E$, its orbit contains
$[q_j]$ and increases the rank. Thus at most $d$ selections suffice, and
full rank proves~\eqref{eq:generators}. Products with a vector-space basis
of the whole algebra $E$ span each orbit, so repeated composition is
unnecessary. The random candidates affect the size of the selected family,
not correctness; minimum cardinality is not asserted. The two action
tables require $2dt$ lift products and quotient-coordinate solves. This
preparation is useful when its cost is outweighed by the image and
preimage computations saved across the pruning passes.

\subsection{Dependencies on the current submodule}
A further reduction applies to any finite candidate family $B\subseteq H$,
including the composition generators. At the start of an I or K pass,
respectively, form the restriction maps
\begin{align*}
 T_I&:\Span(B)\to\Hom(U,N/\eta U),\qquad f\mapsto[f|_U],\\
 T_K&:\Span(B)\to\Hom(V,N/\eta V),\qquad f\mapsto[f|_V].
\end{align*}
Choose subsets $B_I',B_K'\subseteq B$ whose images are bases of the
respective restriction images.

\begin{proposition}
These subsets preserve the I constraints and K image sum:
\begin{align*}
 U\cap\bigcap_{f\in B}f^{-1}(\eta U)
 &=U\cap\bigcap_{s\in B_I'}s^{-1}(\eta U),\\
 \eta V+\sum_{f\in B}f(V)&=\eta V+\sum_{t\in B_K'}t(V).
\end{align*}
No stability under $E$ is needed for these equalities.
\end{proposition}
\begin{proof}
For I, every omitted $f$ satisfies
$(f-\sum_s c_ss)(U)\subseteq\eta U$ for some scalars $c_s$.
Thus the retained constraints imply every omitted constraint.
For K, the corresponding dependence gives
$(f-\sum_t c_tt)(V)\subseteq\eta V$, so every omitted image lies
in the retained sum. The reverse inclusions follow from selecting subsets.
\end{proof}

To compute the restrictions, let $G$ be a generator map for the submodule
at the start of the pass, let $S$ lift generators of its image under
$\eta$, and let $R_N$ be the target relation matrix. For each generator
grade $g$, reduce the corresponding column
of $fG$ modulo the columns of $[\,R_N\mid S\,]$ born at or before $g$.
Cache one reducer per distinct grade, flatten the residual matrices, and
retain independent vectors. Each fixed reducer gives a linear normal form
that vanishes exactly on the target submodule at that grade. Since a map
vanishes if and only if it vanishes on generators, these residuals have
exactly the dependencies of the restricted maps. No presentation of the
current submodule is needed.

The subsets are recomputed for each pass. During an I pass, the working
submodule stays inside the original $U$ and the target $\eta U$ remains
fixed, so the dependencies remain valid as the intersections proceed.
During a K pass, $V$ remains fixed. An empty retained family still gives
the K update $\eta^{-1}(\eta V)$, which includes $\ker\eta$ when $V=0$.
This must not be replaced by an early return.\footnote{The implementation
separately preserves the pre-existing convention of returning $(M,0)$
when the original lift quotient $L_\delta/\tau L_0$ is zero. That
condition is stronger than $Q=0$; an empty reduced family does not
trigger it.}
\end{document}
